\documentclass[11pt,a4paper]{article}
\usepackage[utf8]{inputenc}
\usepackage[T1]{fontenc}
\usepackage[french]{babel}
\usepackage{amsmath,amssymb}
\usepackage{graphicx}
\usepackage[dvipsnames]{xcolor}
\usepackage{float}
\usepackage{booktabs}
\usepackage[margin=1.8cm,top=2cm]{geometry}
\usepackage{hyperref}

\graphicspath{{../test_lagrange_gresho_cluster/}}
\newcommand{\lagdir}{../test_lagrange_gresho_cluster/}
\newcommand{\sch}{vp\_pp}
\newcommand{\schpath}{vp_pp}

\title{Schéma lagrangien \texttt{\sch} --- tourbillon de Gresho}
\author{Vincent Delmas}
\date{\today}

\begin{document}
\maketitle

\noindent
Solveur \texttt{subfvlagrange}, ordre 1, tourbillon de Gresho à $\mathrm{Ma}=10^{-2}$
sur $[-0{,}5;0{,}5]^2$ (1600 hexaèdres / 3200 prismes), parois sur tout le bord.

\medskip
\noindent\textbf{Pourquoi ce cas ouvre l'étude.} En Lagrangien pur il n'y a pas de terme
d'advection : le maillage suit le fluide et seul le couplage nodal pression/vitesse
agit. Le tourbillon de Gresho étant une solution \emph{stationnaire}, tout écart à
l'état initial est une erreur du schéma, sans rien à démêler d'un transport numérique.
C'est donc le test le plus direct du solveur nodal, avant d'ajouter l'advection et de
passer aux schémas Euler.

\medskip
\noindent\textbf{Les trois variantes comparées.}
\begin{itemize}\itemsep2pt
  \item \texttt{classic} : vitesse nodale seule (famille EUCCLHYD, Maire \emph{et al.}), fermeture par
        les pressions unilatérales $p_l-\lambda_l(u_p\!\cdot\!n-v_l)$ ;
  \item \texttt{vp\_pp} : vitesse et pression nodales issues \emph{toutes deux de
        systèmes nodaux}, donc sans aucune longueur ad hoc. Une seule pression sert des
        deux côtés d'une sous-face, ce qui rend le flux symétrique et conservatif.
\end{itemize}

\noindent
\texttt{multi\_point\_pressure} n'est pas transposable ici : une vitesse de n\oe{}ud est
indispensable pour déplacer le maillage.

\section*{Résultats}

\begin{table}[H]\centering
\begin{tabular}{lcc}
\toprule
énergie cinétique restante $E_{kin}/E_{kin,0}$ à $t=0{,}2$ & quad & tri \\
\midrule
\texttt{classic}   & $4\cdot10^{-5}$ & $0$ \\
\texttt{vp\_pp} & $1{,}012$ & $1{,}151$ \\
\bottomrule
\end{tabular}
\end{table}

\noindent
\texttt{classic} dissipe $99{,}996\,\%$ de l'énergie cinétique. La dissipation y opère à
l'échelle \emph{acoustique} et non convective, donc $100\times$ trop vite à
$\mathrm{Ma}=10^{-2}$ : à $\mathrm{Ma}=1$ le même cas n'en perd que $31\,\%$, à
$\mathrm{Ma}=10^{-1}$ $82\,\%$, à $\mathrm{Ma}=10^{-2}$ $99{,}8\,\%$ (mesuré).

\medskip
\noindent
\texttt{vp\_pp} conserve l'énergie totale à 6 chiffres. La dérive de l'énergie
cinétique est un transfert depuis l'énergie interne de $7\cdot10^{-7}$ en relatif,
amplifié par le rapport $E_{int}/E_{kin}=2{,}1\cdot10^5$ propre au bas Mach.

\section*{Convergence bas Mach de la fluctuation de densité}

Même protocole que du côté Euler : $t_{final}=10^{-2}$ et
$L_2(\rho)=\sqrt{\sum_i V_i(\rho_i-1)^2}$, de sorte que les chiffres se comparent
directement à ceux de \texttt{ns\_gresho\_convergence\_main.F90}.

\begin{table}[H]\centering
\begin{tabular}{lcccc}
\toprule
$L_2(\rho)$ & \multicolumn{2}{c}{\texttt{classic}} & \multicolumn{2}{c}{\texttt{vp\_pp}} \\
 & quad & tri & quad & tri \\
\midrule
$\mathrm{Ma}=10^{-1}$ & 3{,}27e-5 & 9{,}28e-5 & 2{,}16e-5 & 1{,}50e-4 \\
$\mathrm{Ma}=10^{-2}$ & 1{,}75e-6 & 3{,}23e-6 & 4{,}07e-7 & 1{,}53e-5 \\
$\mathrm{Ma}=10^{-3}$ & 2{,}25e-8 & 2{,}16e-8 & 3{,}89e-8 & 2{,}47e-6 \\
$\mathrm{Ma}=10^{-4}$ & \textbf{2{,}31e-10} & \textbf{2{,}17e-10} & 3{,}89e-9 & 1{,}94e-6 \\
\midrule
pente & $\approx2$ & $\approx2$ & $\approx1$ & plancher \\
\bottomrule
\end{tabular}
\end{table}

\paragraph{Ce tableau ne doit pas être lu seul.} \texttt{classic} atteint
$\mathrm{Ma}^2$, ce qui en ferait le meilleur schéma du lot. Or l'énergie cinétique
restante \emph{à l'instant même de la mesure} vaut :

\begin{table}[H]\centering
\begin{tabular}{lcc}
\toprule
$E_{kin}(t{=}10^{-2})/E_{kin}(0)$ & $\mathrm{Ma}=10^{-2}$ & $\mathrm{Ma}=10^{-4}$ \\
\midrule
\texttt{classic} quad   & $0{,}226$ & \textcolor{BrickRed}{$\mathbf{0{,}000}$} \\
\texttt{classic} tri    & $0{,}118$ & \textcolor{BrickRed}{$\mathbf{0{,}000}$} \\
\texttt{vp\_pp} quad & $1{,}0000$ & $1{,}0000$ \\
\texttt{vp\_pp} tri  & $1{,}0001$ & $1{,}0009$ \\
\bottomrule
\end{tabular}
\end{table}

À $\mathrm{Ma}=10^{-4}$, \texttt{classic} a \textbf{intégralement détruit le tourbillon
avant la mesure} : sa fluctuation de densité tend vers zéro parce qu'il n'y a plus
d'écoulement, pas parce que le schéma est précis. Le phénomène s'aggrave d'ailleurs
quand le Mach baisse ($22{,}6\,\%$ d'énergie restante à $10^{-2}$, $0\,\%$ à
$10^{-4}$), ce qui produit mécaniquement la pente $\mathrm{Ma}^2$ observée.

\texttt{vp\_pp} conserve exactement l'énergie et se trouve pénalisé par la métrique.
Le critère de convergence en Mach, pris isolément, classe donc les schémas à l'envers ---
constat déjà obtenu côté Euler, et retrouvé ici dans un solveur \emph{sans aucun terme
d'advection}, donc sans rien d'autre à incriminer que le couplage nodal.

\clearpage
\newcommand{\minifig}[2]{%
  \begin{minipage}[t]{0.48\linewidth}\centering
    \IfFileExists{\lagdir#1}%
      {\includegraphics[width=\linewidth]{#1}}%
      {\fbox{\parbox[c][0.32\linewidth]{0.8\linewidth}{\centering\small non disponible\\(le schéma s'est arrêté)}}}%
    \\[2pt]{\small #2}
  \end{minipage}\hfill}

\section*{Maillage --- quad}
\begin{figure}[H]\centering
  \minifig{\schpath_gresho_quad/wireframe_init.png}{maillage initial}
  \minifig{\schpath_gresho_quad/wireframe_final.png}{maillage final}
\end{figure}
\section*{Vitesse --- quad}
\begin{figure}[H]\centering
  \minifig{\schpath_gresho_quad/velocity_init.png}{$|v|$ initial}
  \minifig{\schpath_gresho_quad/velocity_final.png}{$|v|$ final}
\end{figure}

\clearpage
\section*{Maillage --- tri}
\begin{figure}[H]\centering
  \minifig{\schpath_gresho_tri/wireframe_init.png}{maillage initial}
  \minifig{\schpath_gresho_tri/wireframe_final.png}{maillage final}
\end{figure}
\section*{Vitesse --- tri}
\begin{figure}[H]\centering
  \minifig{\schpath_gresho_tri/velocity_init.png}{$|v|$ initial}
  \minifig{\schpath_gresho_tri/velocity_final.png}{$|v|$ final}
\end{figure}

\clearpage
\section*{Maillage --- Voronoï}
Maillage polygonal irrégulier de 1600 cellules --- même nombre que le quad
$40\times40$, donc à résolution égale --- de 4 à 8 côtés (5,70 en moyenne).
Les germes sont construits par réflexion à travers les deux axes médians, si
bien que ce maillage est \emph{miroir-symétrique} par rapport à $x=0$ et $y=0$ :
ce n'est pas un Voronoï générique.
\begin{figure}[H]\centering
  \minifig{\schpath_gresho_voro/wireframe_init.png}{maillage initial}
  \minifig{\schpath_gresho_voro/wireframe_final.png}{maillage final}
\end{figure}
\section*{Vitesse --- Voronoï}
\begin{figure}[H]\centering
  \minifig{\schpath_gresho_voro/velocity_init.png}{$|v|$ initial}
  \minifig{\schpath_gresho_voro/velocity_final.png}{$|v|$ final}
\end{figure}

\paragraph{Le Voronoï n'est pas le cas difficile --- les triangles le sont.}
Le pas de temps final à $t=0{,}25$, à nombre de cellules égal, le montre
directement :
\begin{table}[H]\centering
\begin{tabular}{lcccc}
\toprule
$dt$ à $t=0{,}25$ & quad & tri & Voronoï & Voronoï (Lloyd 20) \\
\midrule
\texttt{classic}   & $2{,}416\cdot10^{-5}$ & \textcolor{BrickRed}{$3{,}667\cdot10^{-7}$} & $1{,}343\cdot10^{-5}$ & $1{,}379\cdot10^{-5}$ \\
\texttt{vp\_pp} & $1{,}010\cdot10^{-5}$ & $5{,}011\cdot10^{-6}$ & $4{,}895\cdot10^{-6}$ & $4{,}353\cdot10^{-6}$ \\
\bottomrule
\end{tabular}
\end{table}
\noindent
\texttt{classic} perd un facteur $66$ sur les triangles, mais seulement $1{,}8$
sur le Voronoï. Le maillage polygonal, qu'on attendrait le plus hostile, est en
fait bien mieux supporté que le maillage triangulaire. Passer de Lloyd 3 à
Lloyd 20 ne change presque rien, ce qui confirme que le résultat sur maillage
plus irrégulier n'est pas un artefact de qualité de maillage.

\paragraph{Énergie cinétique retenue et fluctuation de densité.}
Le balayage en Mach jusqu'à $t_{final}=10^{-2}$, mesuré exactement comme côté
Euler, donne :
\begin{table}[H]\centering
\begin{tabular}{lccc}
\toprule
sur Voronoï & $E_{kin}$ à Ma$=10^{-2}$ & $E_{kin}$ à Ma$=10^{-4}$ & pente $L_2(\rho)$ \\
\midrule
\texttt{classic}   & $46{,}8\%$ & \textcolor{BrickRed}{$0{,}00\%$} & $1{,}91$ \\
\texttt{vp\_pp} & $99{,}75\%$ & $\mathbf{98{,}29\%}$ & $0{,}97$ \\
\bottomrule
\end{tabular}
\end{table}
\noindent
La pente de $1{,}91$ de \texttt{classic}, proche de Mach$^2$, est un
\emph{artefact} : le tourbillon ayant été annihilé, $\rho\equiv1$ trivialement.
C'est la même inversion de classement que sur quad et tri, et que côté Euler.
Les pentes sont obtenues par moindres carrés sur les quatre nombres de Mach,
comme partout ailleurs dans cette étude ; entre les deux derniers points
seulement elles vaudraient $2{,}22$ et $0{,}93$.

\paragraph{Lecture du maillage final sur Voronoï.} Elle est contre-intuitive.
Le maillage de \texttt{classic} ressort quasi intact : n'ayant plus de vitesse,
il n'a plus rien pour déplacer ses n\oe{}uds. Celui de \texttt{vp\_pp} montre
au contraire une spirale centrale nette, trace des cellules réellement
entraînées par la rotation. Ici, un maillage peu déformé est le symptôme de la
dissipation, pas de la qualité.

\clearpage
\section*{Sedov --- robustesse au choc fort}

Le Gresho teste la capacité à \emph{ne rien faire} ; le Sedov teste la survie à une onde
de souffle. Les deux sondent des propriétés opposées, et c'est ce contraste qui est
instructif. Ordre 1, $t_{max}=1$, \texttt{cfl}=0{,}1, domaine $[-1{,}2;1{,}2]^2$.

\begin{table}[H]\centering
\begin{tabular}{lcc}
\toprule
état à $t_{max}$ & quad & tri \\
\midrule
\texttt{classic} : bornes & $[-1{,}2;1{,}2]^2$ & $[-1{,}2;1{,}2]^2$ \\
\texttt{classic} : mailles NaN & $0/3721$ & $0/36332$ \\
\midrule
\texttt{vp\_pp} : mailles NaN & \textcolor{BrickRed}{$44/3721$} & \textcolor{BrickRed}{$25/36332$} \\
\texttt{vp\_pp} : itérations faites & \textcolor{BrickRed}{$0$} & \textcolor{BrickRed}{$0$} \\
\texttt{vp\_pp} : bornes des n\oe{}uds finis & \textcolor{BrickRed}{$[-6{,}63;6{,}63]^2$} & $[-1{,}2;1{,}2]^2$ \\
\bottomrule
\end{tabular}
\end{table}

\noindent
\texttt{vp\_pp} diverge en moins de $100$ itérations --- le journal n'écrit une
ligne que tous les $100$ pas et n'en contient aucune. Le souffle \emph{a} démarré :
l'instrumentation du schéma montre $\tau = 0{,}990$ et $1{,}013$ au lieu de $1$, et
une pression ambiante montée de $10^{-12}$ à $\sim\!2\cdot10^{-8}$.

\paragraph{La cause est une perte de positivité de l'énergie interne.} Dans le
quasi-vide, l'énergie ambiante initiale ($10^{-12}/(\gamma-1) \approx
2{,}5\cdot10^{-12}$) est négligeable devant le travail reçu ($\sim\!10^{-8}$) : une
maille peut donc fournir plus de travail qu'elle n'a d'énergie, et $\mathrm{sol}(5)$
devient négatif. La pression suit, puis $\lambda = \max(\sqrt{\gamma p \tau}/\tau,
\dots)$ prend la racine d'un argument négatif et vaut NaN. Les deux autres termes du
$\max$ sont protégés, le terme acoustique ne l'est pas. $M_p$ est donc NaN
\emph{avant} l'inversion : le solveur nodal n'est pas en cause, c'est l'accumulation.

\paragraph{Pourquoi \texttt{vp\_pp} et pas \texttt{classic}.} \texttt{vp\_pp} ferme
avec une pression nodale \emph{unique} des deux côtés d'une sous-face : le flux est
conservatif, mais rien ne garantit la positivité maille par maille. \texttt{classic}
ferme au contraire avec des pressions unilatérales $p_l-\lambda_l(u_p\!\cdot\!n-v_l)$,
cohérentes avec l'état local.

\paragraph{Attention à ce que montre le fichier de sortie.} Il est écrit \emph{après}
que $t$ soit devenu NaN, et ne reflète donc pas l'état au moment de la divergence. Sur
quad il porte $44$ mailles NaN sur $3721$, une pression négative résiduelle
($-1{,}036\cdot10^{-3}$) et huit n\oe{}uds projetés jusqu'à $5{,}5\times$ la taille du
domaine ; sur tri, $25$ mailles NaN, aucune pression négative visible et des bornes
intactes. Cette différence est \textbf{trompeuse} : la pression négative existe aussi
sur tri, mais transitoire, dans une maille qui finit NaN dans le fichier. Sur ces
calculs, ni les bornes du maillage ni les champs finaux ne sont un diagnostic fiable
--- seul le comptage de NaN l'est.

\paragraph{Attention : ce run se termine avec \texttt{EXIT\_CODE: 0}.} Une fois $dt$
devenu NaN, la boucle \texttt{do while (t < t\_max)} est fausse au premier test et le
code sort proprement, en écrivant \texttt{run.timestamp} et trois VTU. Sans inspecter
les champs ou les bornes du maillage, ce calcul compte pour une réussite.

\paragraph{Pourquoi cet échec était prévisible.} \texttt{vp\_pp} utilise une pression
nodale \emph{unique} des deux côtés d'une sous-face : le flux est symétrique, donc
conservatif, mais sans production d'entropie. C'est exactement ce qui lui permet de
conserver l'énergie du tourbillon de Gresho --- et exactement ce qui lui manque pour
encaisser un choc fort. \texttt{classic} ferme au contraire avec des pressions
\emph{unilatérales} $p_l-\lambda_l(u_p\!\cdot\!n-v_l)$, dont l'asymétrie \emph{est} le
mécanisme dissipatif : inutile en bas Mach, indispensable ici.

La même dichotomie s'observe côté Euler, dans un formalisme pourtant différent :
\texttt{multi\_point\_pressure} (pression nodale) est le meilleur sur le flux de chaleur
quad et excellent en bas Mach, mais s'arrête sur \texttt{sedov\_tri}, \texttt{sedov\_hex}
et le demi-cylindre tri ; \texttt{multi\_point} (vitesse nodale) est robuste partout et
dissipatif partout.

\begin{figure}[H]\centering
  \minifig{../test_lagrange_sedov_cluster/\schpath_sedov_2d_quad/wireframe_final.png}{maillage final --- quad}
  \minifig{../test_lagrange_sedov_cluster/\schpath_sedov_2d_quad/density_final.png}{densité --- quad}
\end{figure}
\begin{figure}[H]\centering
  \minifig{../test_lagrange_sedov_cluster/\schpath_sedov_2d_tri/wireframe_final.png}{maillage final --- tri}
  \minifig{../test_lagrange_sedov_cluster/\schpath_sedov_2d_tri/density_final.png}{densité --- tri}
\end{figure}

\vfill
\noindent\textbf{Lecture du maillage final.} Un maillage final peu déformé n'est pas un
bon signe en soi : si le schéma a dissipé le tourbillon, il n'y a plus de vitesse pour
déplacer les n\oe{}uds. C'est exactement ce qui se produit pour \texttt{classic}, dont
le maillage paraît sain alors que la physique a disparu. Les deux figures doivent donc
se lire ensemble.

\end{document}
